\documentclass{article}
\usepackage[T1]{fontenc}
\usepackage{lmodern}
\usepackage{graphicx}
\usepackage{amsmath,amssymb}
\usepackage[a4paper,margin=2cm]{geometry}
\usepackage[absolute,overlay]{textpos}
\setlength{\TPHorizModule}{1mm}
\setlength{\TPVertModule}{1mm}
\usepackage[indonesian]{babel}
\usepackage{hyperref}
\setlength{\emergencystretch}{2em}
% unit: Halaman 14

\begin{document}

% === Halaman 14 (llm) ===

\section*{Metode Steganalisis}

1. Serangan berbasis visual (\textit{visual attacks})
\begin{itemize}
    \item Khusus untuk \textit{stego-object} berupa citra
    \item Bersifat subjektif, karena melakukan pengamatan secara kasat mata dengan melihat artefak yang mencurigakan di dalam \textit{stego-image}, lalu membandingkannya dengan citra asli (\textit{cover image})
    \item Digunakan pada masa-masa awal riset steganalisis
    \item Contoh serangan visual:
    \begin{enumerate}
        \item \textit{LSB plane attack}
        \item \textit{Filtered visual attack} (Enhanced LSB)
    \end{enumerate}
\end{itemize}

\end{document}
